%%%%%%%%%%%%%%%%
%%% exod 31 %%%
%%%%%%%%%%%%%%%%

\Chapter{31}

\Verse{1}
{
	\hP{וַיְדַבֵּר יְהֹוָה אֶל מֹשֶׁה לֵּאמֹר׃}
}
{
	\eP{And YHWH spoke to Moses, saying,}
}
{
%	\footnote{
%		Bezalel. Bezalel is the first artist, and he is summoned by name. The
%		irony is that immediately after this a work of art, the golden calf, is
%		made! The special importance of the artist and his or her relationship
%		to the work of art is conveyed as Aaron makes the calf (32:4) but then,
%		when he has to explain it to Moses, speaks as if it somehow came to
%		exist by itself: “I threw it into the fire, and this calf came out!”
%		Some have taken the name Bezalel to mean “in the shadow of God” (bl ’l).
%		To me it also intimates “in the image of God” (belem ’lhîm) from the
%		creation story (Gen 1:27). (On the absence of the letter mem, compare
%		the name Noah being related to the root nm, and the name Samuel being
%		related to the root ’l. See the comment on Gen 5:29.) And Bezalel is
%		filled with the “spirit of God,” which also comes from the creation
%		story (1:2). The allusions to creation are attractive because Bezalel,
%		after all, as the great artist of the Torah, is the creative one, who
%		fashions the Tabernacle and its contents, including the ark. Being
%		creative is the ultimate imitatio Dei. And the conclusion of the account
%		of the construction of the Tabernacle uses some of the same language as
%		the conclusion of the creation account. (See the comment on Exod 39:32.)
%	}
}

\Verse{2}
{
	\hP{רְאֵה קָרָאתִי בְשֵׁם בְּצַלְאֵל בֶּן אוּרִי בֶן חוּר לְמַטֵּה יְהוּדָה׃}
}
{
	\eP{
		“See, I’ve called by name Bezalel, son of Uri, son of Hur,
		of the tribe of Judah.
	}
}
{
}

\Verse{3}
{
	\hP{וָאֲמַלֵּא אֹתוֹ רוּחַ אֱלֹהִים בְּחכְמָה וּבִתְבוּנָה וּבְדַעַת וּבְכל מְלָאכָה׃}
}
{
	\eP{
		And I’ve filled him with the spirit of {\God} in wisdom and
		in understanding and in knowledge and in every kind of
		work,
	}
}
{
}

\Verse{4}
{
	\hP{לַחְשֹׁב מַחֲשָׁבֹת לַעֲשׂוֹת בַּזָּהָב וּבַכֶּסֶף וּבַנְּחֹשֶׁת׃}
}
{
	\eP{
		to form conceptions to make in gold and in silver and in
		bronze
	}
}
{
}

\Verse{5}
{
	\hP{וּבַחֲרֹשֶׁת אֶבֶן לְמַלֹּאת וּבַחֲרֹשֶׁת עֵץ לַעֲשׂוֹת בְּכל מְלָאכָה׃}
}
{
	\eP{
		and in cutting stone for setting and in cutting wood—for
		making things in every kind of work.
	}
}
{
}

\Verse{6}
{
	\hP{וַאֲנִי הִנֵּה נָתַתִּי אִתּוֹ אֵת אהֳלִיאָב בֶּן אֲחִיסָמָךְ לְמַטֵּה דָן וּבְלֵב כּל חֲכַם לֵב נָתַתִּי חכְמָה וְעָשׂוּ אֵת כּל אֲשֶׁר צִוִּיתִךָ׃}
}
{
	\eP{
		And I: here, I’ve put Oholiab, son of Ahisamach, of the
		tribe of Dan, with him. And in the heart of every
		wise-hearted person I’ve put wisdom. And they’ll make
		everything that I’ve commanded you:
	}
}
{
}

\Verse{7}
{
	\hP{אֵת אֹהֶל מוֹעֵד וְאֶת הָאָרֹן לָעֵדֻת וְאֶת הַכַּפֹּרֶת אֲשֶׁר עָלָיו וְאֵת כּל כְּלֵי הָאֹהֶל׃}
}
{
	\eP{
		the Tent of Meeting and the Ark of the Testimony and the
		atonement dais that is on it and all the equipment of the
		tent
	}
}
{
}

\Verse{8}
{
	\hP{וְאֶת הַשֻּׁלְחָן וְאֶת כֵּלָיו וְאֶת הַמְּנֹרָה הַטְּהֹרָה וְאֶת כּל כֵּלֶיהָ וְאֵת מִזְבַּח הַקְּטֹרֶת׃}
}
{
	\eP{
		and the table and its equipment and the pure menorah and
		all its equipment and the incense altar
	}
}
{
}

\Verse{9}
{
	\hP{וְאֶת מִזְבַּח הָעֹלָה וְאֶת כּל כֵּלָיו וְאֶת הַכִּיּוֹר וְאֶת כַּנּוֹ׃}
}
{
	\eP{
		and the burnt offering altar and all its equipment and the
		basin and its stand
	}
}
{
}

\Verse{10}
{
	\hP{וְאֵת בִּגְדֵי הַשְּׂרָד וְאֶת בִּגְדֵי הַקֹּדֶשׁ לְאַהֲרֹן הַכֹּהֵן וְאֶת בִּגְדֵי בָנָיו לְכַהֵן׃}
}
{
	\eP{
		and the fabric clothes and the holy clothes for Aaron, the
		priest, and his sons’ clothes for functioning as priests
	}
}
{
}

\Verse{11}
{
	\hP{וְאֵת שֶׁמֶן הַמִּשְׁחָה וְאֶת קְטֹרֶת הַסַּמִּים לַקֹּדֶשׁ כְּכֹל אֲשֶׁר צִוִּיתִךָ יַעֲשׂוּ׃}
}
{
	\eP{
		and the anointing oil and the incense of fragrances for
		the Holy. According to everything that I’ve commanded you,
		they shall do.”
	}
}
{
}

\Verse{12}
{
	\hP{וַיֹּאמֶר יְהֹוָה אֶל מֹשֶׁה לֵּאמֹר׃}
}
{
	\eP{And YHWH said to Moses, saying,}
}
{
}

\Verse{13}
{
	\hP{וְאַתָּה דַּבֵּר אֶל בְּנֵי יִשְׂרָאֵל לֵאמֹר אַךְ אֶת שַׁבְּתֹתַי תִּשְׁמֹרוּ כִּי אוֹת הִוא בֵּינִי וּבֵינֵיכֶם לְדֹרֹתֵיכֶם לָדַעַת כִּי אֲנִי יְהֹוָה מְקַדִּשְׁכֶם׃}
}
{
	\eP{
		“And you, speak to the children of Israel, saying, ‘Just:
		you shall observe my Sabbaths, because it is a sign
		between me and you through your generations: to know that
		I, YHWH, make you holy.
	}
}
{
%	\footnote{
%		Just. The force of this word here seems to be to say that the section
%		preceding this tells what the skilled craftsmen will do, and what God
%		will do, to produce the holy place and all the holy objects; and now it
%		will be specified what the people must do to become holy themselves. Why
%		just the Sabbath out of all the commandments? Because “it is a sign
%		between me and you.” As the sign of the Noahic covenant is the rainbow
%		and the sign of the Abrahamic covenant is circumcision, so the sign of
%		the Israelite covenant is the Sabbath. A sign participates in the thing
%		that it symbolizes. (Thus, at first a flag is just a country’s symbol,
%		but later people will say, “I’d die for the flag!”) And so its
%		fulfillment is of exceptional importance. The rainbow sign is provided
%		by God, and so it is certain. The signs of the other two covenants must
%		be performed by humans, and so they are commanded with extreme emphasis
%		and strictness. Both declare that failure to perform the sign will
%		result in “cutting off” (Gen 17:14; Exod 31:14).
%	}
}

\Verse{14}
{
	\hP{וּשְׁמַרְתֶּם אֶת הַשַּׁבָּת כִּי קֹדֶשׁ הִוא לָכֶם מְחַלְלֶיהָ מוֹת יוּמָת כִּי כּל הָעֹשֶׂה בָהּ מְלָאכָה וְנִכְרְתָה הַנֶּפֶשׁ הַהִוא מִקֶּרֶב עַמֶּיהָ׃}
}
{
	\eP{
		And you shall observe the Sabbath, because it is a holy
		thing to you. One who desecrates it shall be put to death.
		Because anyone who does work in it: that person will be
		cut off from among his people. 31:14,15. put to death. The
		extreme punishment of execution for violation of the
		Sabbath must be understood in the context of the singular
		importance of the Sabbath: as sign of the covenant and as
		imitatio Dei—see the previous comment. Also, as I
		discussed in my comments on Genesis 1, the special place
		of the Sabbath in the creation of the universe marks it as
		having cosmic significance as well as historical
		significance. Combining the cosmic with history, it sets
		the history of human events that begins with Genesis 1 in
		the context of the very nature of the universe. And it
		makes time holy. Thus this long list of holy objects
		concludes with this reminder that periods of time are holy
		as well. The Sabbath is thus so phenomenally important
		that a violation of the Sabbath bears an extraordinary
		punishment.
	}
}
{
}

\Verse{15}
{
	\hP{שֵׁשֶׁת יָמִים יֵעָשֶׂה מְלָאכָה וּבַיּוֹם הַשְּׁבִיעִי שַׁבַּת שַׁבָּתוֹן קֹדֶשׁ לַיהֹוָה כּל הָעֹשֶׂה מְלָאכָה בְּיוֹם הַשַּׁבָּת מוֹת יוּמָת׃}
}
{
	\eP{
		Six days work shall be done, and in the seventh day is a
		Sabbath, a ceasing, a holy thing to YHWH. Anyone who does
		work in the Sabbath day shall be put to death.
	}
}
{
}

\Verse{16}
{
	\hP{וְשָׁמְרוּ בְנֵי יִשְׂרָאֵל אֶת הַשַּׁבָּת לַעֲשׂוֹת אֶת הַשַּׁבָּת לְדֹרֹתָם בְּרִית עוֹלָם׃}
}
{
	\eP{
		And the children of Israel shall observe the Sabbath, to
		make the Sabbath through their generations, an eternal
		covenant.
	}
}
{
}

\Verse{17}
{
	\hP{בֵּינִי וּבֵין בְּנֵי יִשְׂרָאֵל אוֹת הִוא לְעֹלָם כִּי שֵׁשֶׁת יָמִים עָשָׂה יְהֹוָה אֶת הַשָּׁמַיִם וְאֶת הָאָרֶץ וּבַיּוֹם הַשְּׁבִיעִי שָׁבַת וַיִּנָּפַשׁ׃}
}
{
	\eP{
		Between me and the children of Israel it is a sign
		forever, because for six days YHWH made the skies and the
		earth, and in the seventh day He ceased and was
		refreshed.’”
	}
}
{
%	\footnote{
%		refreshed. This verb occurs only twice in the Torah: here, referring to
%		God’s Sabbath rest; and in Exod 23:12, where it refers to the Sabbath
%		rest of “your maid’s son and the alien.” It is remarkable that the
%		parallel cases are the son of a slave and an alien. And this comes in
%		the middle of one of the most blatant cases of imitatio Dei in the
%		Torah: humans are to cease on the Sabbath because God ceased creating on
%		the Sabbath. This now conveys that every human participates in this
%		divine phenomenon. One cannot rest but still require his or her slave to
%		work. One cannot rest but still count on non-Israelites in the community
%		to do the work.
%	}
}

\Verse{18}
{
	\hP{וַיִּתֵּן אֶל מֹשֶׁה כְּכַלֹּתוֹ לְדַבֵּר אִתּוֹ בְּהַר סִינַי שְׁנֵי לֻחֹת הָעֵדֻת לֻחֹת אֶבֶן כְּתֻבִים בְּאֶצְבַּע אֱלֹהִים׃}
}
{
	\eP{
		And when He finished speaking with him in Mount Sinai, He
		gave the two tablets of the Testimony to Moses, tablets of
		stone, written by the finger of {\God}.
	}
}
{
%	\footnote{
%		two tablets of the Testimony. There are two tablets of
%		testimony/witness, and there must always be at least two witnesses to
%		give testimony in law. And Moses names two witnesses, heaven and earth,
%		concerning the people’s future at the end of the Torah (Deut 31:28;
%		32:1).
%	}
}
